\begin{ex}
Phân số đảo ngược của $\frac{3}{5}$ là phân số nào?
\shortans{5/3}
\loigiai{
Là $\frac{5}{3}$.
}
\end{ex}

\begin{ex}
Phân số đảo ngược của $\frac{1}{4}$ là số tự nhiên nào?
\shortans{4}
\loigiai{
$\frac{4}{1} = 4$.
}
\end{ex}

\begin{ex}
Phân số đảo ngược của $\frac{7}{2}$ là phân số nào?
\shortans{2/7}
\loigiai{
Là $\frac{2}{7}$.
}
\end{ex}

\begin{ex}
Tính: $\frac{1}{2} : \frac{1}{3} = \dots$
\shortans{3/2}
\loigiai{
$\frac{1}{2} \times \frac{3}{1} = \frac{3}{2}$.
}
\end{ex}

\begin{ex}
Tính: $\frac{3}{4} : \frac{1}{2} = \dots$
\shortans{3/2}
\loigiai{
$\frac{3}{4} \times \frac{2}{1} = \frac{6}{4} = \frac{3}{2}$.
}
\end{ex}

\begin{ex}
Tính: $\frac{2}{5} : \frac{3}{4} = \dots$
\shortans{8/15}
\loigiai{
$\frac{2}{5} \times \frac{4}{3} = \frac{8}{15}$.
}
\end{ex}

\begin{ex}
Tính: $\frac{4}{7} : \frac{4}{7} = \dots$
\shortans{1}
\loigiai{
Số nào chia cho chính nó cũng bằng 1.
}
\end{ex}

\begin{ex}
Tính: $\frac{5}{8} : \frac{5}{6} = \dots$ (tối giản).
\shortans{3/4}
\loigiai{
$\frac{5}{8} \times \frac{6}{5} = \frac{6}{8} = \frac{3}{4}$.
}
\end{ex}

\begin{ex}
Tính: $\frac{2}{3} : 4 = \dots$ (tối giản).
\shortans{1/6}
\loigiai{
$\frac{2}{3 \times 4} = \frac{2}{12} = \frac{1}{6}$.
}
\end{ex}

\begin{ex}
Tính: $\frac{3}{5} : 6 = \dots$ (tối giản).
\shortans{1/10}
\loigiai{
$\frac{3}{5 \times 6} = \frac{3}{30} = \frac{1}{10}$.
}
\end{ex}

\begin{ex}
Tính: $6 : \frac{2}{3} = \dots$
\shortans{9}
\loigiai{
$6 \times \frac{3}{2} = 9$.
}
\end{ex}

\begin{ex}
Tính: $8 : \frac{4}{5} = \dots$
\shortans{10}
\loigiai{
$8 \times \frac{5}{4} = 10$.
}
\end{ex}

\begin{ex}
Tính: $10 : \frac{5}{2} = \dots$
\shortans{4}
\loigiai{
$10 \times \frac{2}{5} = 4$.
}
\end{ex}

\begin{ex}
Một hình chữ nhật có diện tích $\frac{2}{3}\text{ m}^2$, chiều rộng là $\frac{1}{2}$ m. Chiều dài hình chữ nhật là bao nhiêu mét?
\shortans{4/3}
\loigiai{
$\frac{2}{3} : \frac{1}{2} = \frac{4}{3}$ m.
}
\end{ex}

\begin{ex}
Một hình chữ nhật có diện tích $\frac{3}{5}\text{ m}^2$, chiều dài là $\frac{3}{4}$ m. Chiều rộng hình chữ nhật là bao nhiêu mét?
\shortans{4/5}
\loigiai{
$\frac{3}{5} : \frac{3}{4} = \frac{4}{5}$ m.
}
\end{ex}

\begin{ex}
Một sợi dây dài $\frac{8}{9}$ m được cắt thành 4 đoạn bằng nhau. Mỗi đoạn dài bao nhiêu mét? (tối giản).
\shortans{2/9}
\loigiai{
$\frac{8}{9} : 4 = \frac{2}{9}$ m.
}
\end{ex}

\begin{ex}
Rót $\frac{3}{4}$ lít mật ong vào các chai, mỗi chai $\frac{1}{4}$ lít. Rót được bao nhiêu chai?
\shortans{3}
\loigiai{
$\frac{3}{4} : \frac{1}{4} = 3$ chai.
}
\end{ex}

\begin{ex}
Tìm $x$, biết: $x \times \frac{3}{5} = \frac{6}{7}$.
\shortans{10/7}
\loigiai{
$x = \frac{6}{7} : \frac{3}{5} = \frac{6}{7} \times \frac{5}{3} = \frac{10}{7}$.
}
\end{ex}

\begin{ex}
Tìm $x$, biết: $\frac{2}{3} \times x = 4$.
\shortans{6}
\loigiai{
$x = 4 : \frac{2}{3} = 4 \times \frac{3}{2} = 6$.
}
\end{ex}

\begin{ex}
Tìm $x$, biết: $x : \frac{4}{5} = 5$.
\shortans{4}
\loigiai{
$x = 5 \times \frac{4}{5} = 4$.
}
\end{ex}

\begin{ex}
Tính: $\frac{7}{10} : \frac{7}{10} = \dots$
\shortans{1}
\loigiai{
1.
}
\end{ex}

\begin{ex}
Tính: $\frac{3}{8} : \frac{3}{4} = \dots$ (tối giản).
\shortans{1/2}
\loigiai{
$\frac{3}{8} \times \frac{4}{3} = \frac{1}{2}$.
}
\end{ex}

\begin{ex}
Tính: $\frac{5}{12} : \frac{5}{6} = \dots$ (tối giản).
\shortans{1/2}
\loigiai{
$\frac{5}{12} \times \frac{6}{5} = \frac{1}{2}$.
}
\end{ex}

\begin{ex}
Có 12 kg đường chia đều vào các túi, mỗi túi $\frac{3}{4}$ kg. Chia được bao nhiêu túi?
\shortans{16}
\loigiai{
$12 : \frac{3}{4} = 12 \times \frac{4}{3} = 16$ túi.
}
\end{ex}

\begin{ex}
Tính: $0 : \frac{5}{8} = \dots$
\shortans{0}
\loigiai{
0.
}
\end{ex}
